% LuaLaTeX 컴파일 필요
\documentclass[aspectratio=169]{beamer}

% === 패키지 ===
\usepackage{fontspec}
\usepackage{kotex}
\usepackage{graphicx}
\usepackage{xcolor}
\usepackage{emoji}
\usepackage{amsmath}
\usepackage{booktabs}
\usepackage{array}
\usepackage{colortbl}
\usepackage{tikz}

% === 테마 (metropolis) ===
\usetheme[progressbar=frametitle]{metropolis}

% === 폰트 (한글 + 이모지) ===
\setmainfont{Noto Sans CJK KR}
\setsansfont{Noto Sans CJK KR}
\setmonofont{Noto Sans Mono CJK KR}[Scale=0.85]
\setemojifont{Noto Color Emoji}

% === 색상 (프로젝트 팔레트) ===
\definecolor{primary}{HTML}{1976D2}
\definecolor{darkblue}{HTML}{1565C0}
\definecolor{teal}{HTML}{00897B}
\definecolor{amber}{HTML}{FFA000}
\definecolor{lightblue}{HTML}{E3F2FD}
\definecolor{tableheader}{HTML}{BBDEFB}
\definecolor{lightteal}{HTML}{E0F2F1}

\setbeamercolor{frametitle}{bg=primary, fg=white}
\setbeamercolor{progress bar}{fg=primary}
\setbeamercolor{title separator}{fg=primary}
\setbeamercolor{alerted text}{fg=primary}
\setbeamercolor{example text}{fg=teal}
\setbeamercolor{block title}{bg=primary, fg=white}
\setbeamercolor{block body}{bg=lightblue}

% === 슬라이드 번호 ===
\setbeamertemplate{frame numbering}[fraction]

% === 표 행 간격 ===
\renewcommand{\arraystretch}{1.3}

% === 문서 정보 ===
\title{기계학습의 원리}
\subtitle{AI디지털리터러시 — 1주차 2차시}
\author{청주교육대학교}
\date{2026학년도}

\begin{document}

% === 타이틀 ===
\begin{frame}
\titlepage
\end{frame}

% === 목차 ===
\begin{frame}{오늘의 수업 흐름}
\begin{center}
\begin{tabular}{cl}
\emoji{clapper-board} & \textbf{도입} — AI, ML, DL의 포함 관계 (5분) \\[0.3cm]
\emoji{blue-book} & \textbf{전개 1} — 기계학습의 정의 (5분) \\[0.3cm]
\emoji{green-book} & \textbf{전개 2} — 머신러닝 3가지 방법 + AI 사과농장 시뮬레이터 (30분) \\[0.3cm]
\emoji{orange-book} & \textbf{전개 3} — 머신러닝 학습 과정 4단계 (5분) \\[0.3cm]
\emoji{memo} & \textbf{정리} — 핵심 요약 + 엔트리 연결 (5분) \\
\end{tabular}
\end{center}
\end{frame}

%%%%%%%%%%%%%%%%%%%%%%%%%%%%%%%%%%%%%%%%%%%%%
\section{\emoji{clapper-board} 도입: AI, ML, DL의 포함 관계}
%%%%%%%%%%%%%%%%%%%%%%%%%%%%%%%%%%%%%%%%%%%%%

% --- 그림 01: AI ⊃ ML ⊃ DL ---
\begin{frame}{AI, ML, DL의 포함 관계}
\begin{center}
\Large ``1차시에서 배운 AI, 그 안에는 무엇이 있을까요?''
\end{center}
\vspace{0.2cm}
\begin{center}
\includegraphics[width=0.55\textwidth]{images/fig01_AI_ML_DL포함관계.png}
\end{center}
\end{frame}

\begin{frame}{AI $\supset$ ML $\supset$ DL}
\begin{columns}
\begin{column}{0.33\textwidth}
\begin{block}{\emoji{robot} AI (인공지능)}
\small \textbf{가장 큰 개념}\\
인간 지능 모방하는\\모든 기술\\[0.2cm]
\footnotesize 예: 체스 프로그램,\\스팸 필터, ChatGPT
\end{block}
\end{column}
\begin{column}{0.33\textwidth}
\begin{block}{\emoji{brain} ML (기계학습)}
\small \textbf{AI의 하위 개념}\\
데이터로부터\\패턴 자동 학습\\[0.2cm]
\footnotesize 예: 추천 시스템,\\날씨 예측
\end{block}
\end{column}
\begin{column}{0.33\textwidth}
\begin{block}{\emoji{fire} DL (딥러닝)}
\small \textbf{ML의 하위 개념}\\
인공신경망을 여러 층으로\\쌓아 복잡한 패턴 학습\\[0.2cm]
\footnotesize 예: ChatGPT, AlphaGo
\end{block}
\end{column}
\end{columns}
\vspace{0.3cm}
\begin{center}
\small \emoji{light-bulb} \textbf{비유}: AI = 모든 차량, ML = 엔진 있는 차량, DL = 첨단 전기차
\end{center}
\end{frame}

%%%%%%%%%%%%%%%%%%%%%%%%%%%%%%%%%%%%%%%%%%%%%
\section{\emoji{blue-book} 전개 1: 기계학습의 정의}
%%%%%%%%%%%%%%%%%%%%%%%%%%%%%%%%%%%%%%%%%%%%%

% --- 그림 02: 전통적 프로그래밍 vs 기계학습 ---
\begin{frame}{전통적 프로그래밍 vs 기계학습}
\begin{center}
\includegraphics[width=0.75\textwidth]{images/fig02_전통적프로그래밍vs기계학습.png}
\end{center}
\end{frame}

\begin{frame}{두 방식의 핵심 차이}
\begin{columns}
\begin{column}{0.48\textwidth}
\begin{block}{\emoji{desktop-computer} 전통적 프로그래밍}
\centering
\textbf{규칙} + 데이터 $\to$ 결과\\[0.3cm]
\small 사람이 규칙을 직접 작성\\
``무료'' + ``당첨'' $\to$ 스팸\\[0.2cm]
\footnotesize \emoji{warning} 규칙을 모두 나열할 수 없음
\end{block}
\end{column}
\begin{column}{0.48\textwidth}
\begin{block}{\emoji{brain} 기계학습}
\centering
데이터 + 결과 $\to$ \textbf{규칙(패턴)}\\[0.3cm]
\small 컴퓨터가 패턴을 스스로 발견\\
수천 개 메일 $\to$ 스팸 패턴\\[0.2cm]
\footnotesize \emoji{check-mark-button} 새로운 패턴도 자동 적응
\end{block}
\end{column}
\end{columns}
\end{frame}

% --- 그림 18: 왜 기계학습이 필요한가 ---
\begin{frame}{왜 기계학습이 필요할까?}
\begin{center}
\includegraphics[width=0.65\textwidth]{images/fig18_왜기계학습이필요한가.jpg}
\end{center}
\vspace{0.1cm}
\begin{center}
\small 규칙으로 해결하기 \textbf{어려운 문제}가 있기 때문!
\end{center}
\end{frame}

%%%%%%%%%%%%%%%%%%%%%%%%%%%%%%%%%%%%%%%%%%%%%
\section{\emoji{green-book} 전개 2: 머신러닝의 3가지 방법}
%%%%%%%%%%%%%%%%%%%%%%%%%%%%%%%%%%%%%%%%%%%%%

% === Part A: 이론 (12분) ===

\begin{frame}{머신러닝의 3가지 학습 유형}
\begin{center}
\begin{tabular}{|>{\centering\arraybackslash}p{3cm}|>{\centering\arraybackslash}p{3.5cm}|>{\centering\arraybackslash}p{3.5cm}|>{\centering\arraybackslash}p{3cm}|}
\hline
\rowcolor{tableheader}
& \textbf{지도학습} & \textbf{비지도학습} & \textbf{강화학습} \\
\hline
\textbf{핵심} & 정답이 있는 데이터 & 정답이 없는 데이터 & 보상으로 학습 \\
\hline
\textbf{비유} & 선생님과 공부 & 혼자서 탐색 & 시행착오 학습 \\
\hline
\textbf{사과농장} & 1$\cdot$2$\cdot$4$\cdot$5장 & 3장 & — \\
\hline
\textbf{대표 기술} & KNN, 선형회귀, 의사결정트리 & K-Means & Q-러닝 \\
\hline
\end{tabular}
\end{center}
\vspace{0.3cm}
\begin{center}
\small \emoji{light-bulb} 이 3가지를 기억해 두세요. 곧 \textbf{사과농장}에서 하나씩 만나게 됩니다!
\end{center}
\end{frame}

% --- 지도학습 ---
\begin{frame}{지도학습 (Supervised Learning)}
\begin{block}{\emoji{teacher} 정의}
\textbf{정답(라벨)}이 있는 데이터로 학습하여, 새로운 데이터의 정답을 예측
\end{block}
\vspace{0.3cm}
\begin{center}
\includegraphics[width=0.7\textwidth]{images/fig03_지도학습vs비지도학습.png}
\end{center}
\end{frame}

% --- 비지도학습 ---
\begin{frame}{비지도학습 (Unsupervised Learning)}
\begin{block}{\emoji{magnifying-glass-tilted-left} 정의}
\textbf{정답 없이} 데이터 속 숨겨진 패턴이나 그룹을 스스로 발견
\end{block}
\vspace{0.2cm}
\begin{columns}
\begin{column}{0.48\textwidth}
\begin{center}
\includegraphics[width=\textwidth]{images/fig04_라벨없는_3가지이유.png}
\end{center}
\end{column}
\begin{column}{0.48\textwidth}
\begin{center}
\includegraphics[width=\textwidth]{images/fig05_군집화_3단계과정.png}
\end{center}
\end{column}
\end{columns}
\end{frame}

% --- 강화학습 ---
\begin{frame}{강화학습 (Reinforcement Learning)}
\begin{block}{\emoji{video-game} 정의}
\textbf{보상(+)}과 \textbf{벌점($-$)}을 통해 최적의 행동을 스스로 학습
\end{block}
\vspace{0.2cm}
\begin{columns}
\begin{column}{0.48\textwidth}
\begin{center}
\includegraphics[width=0.9\textwidth]{images/fig07_강화학습루프.png}
\end{center}
\end{column}
\begin{column}{0.48\textwidth}
\begin{center}
\includegraphics[width=0.9\textwidth]{images/fig08_강아지훈련.png}
\end{center}
\end{column}
\end{columns}
\end{frame}

\begin{frame}{강화학습 — 게임으로 이해하기}
\begin{center}
\includegraphics[width=0.55\textwidth]{images/fig09_게임점수획득.png}
\end{center}
\vspace{0.1cm}
\begin{center}
\small 먹이 획득(+10) vs 벽 충돌($-$100) $\to$ 최적 경로 학습!
\end{center}
\end{frame}

% === Part B: AI 사과농장 이야기 — 시뮬레이터 시연 (18분) ===

\begin{frame}[standout]
\emoji{green-apple} AI 사과농장 이야기

\vspace{0.3cm}
시뮬레이터 시연
\end{frame}

% --- 사과농장 도입 ---
\begin{frame}{AI 사과농장 이야기: 새로운 시작}
\begin{center}
\includegraphics[width=0.65\textwidth]{images/fig11_AI사과농장이야기.png}
\end{center}
\vspace{0.1cm}
\begin{center}
\small \emoji{man-farmer} 농부 김사과와 \emoji{robot} AI 로봇 알파가 함께 사과 농장을 운영합니다!
\end{center}
\end{frame}

\begin{frame}{사과농장 — 5가지 ML 기술}
\begin{center}
\begin{tabular}{|>{\centering\arraybackslash}p{0.8cm}|p{4cm}|>{\centering\arraybackslash}p{2.5cm}|>{\centering\arraybackslash}p{2cm}|}
\hline
\rowcolor{tableheader}
\textbf{장} & \textbf{문제 상황} & \textbf{ML 기술} & \textbf{학습 유형} \\
\hline
1장 & \emoji{deciduous-tree} 익은 사과만 골라 따기 & 딥러닝 CNN & 지도학습 \\
\hline
2장 & \emoji{magnifying-glass-tilted-left} 섞인 품종 구분하기 & KNN & 지도학습 \\
\hline
3장 & \emoji{package} 크기별 등급 나누기 & K-Means & 비지도학습 \\
\hline
4장 & \emoji{money-bag} 품질로 가격 매기기 & 선형 회귀 & 지도학습 \\
\hline
5장 & \emoji{delivery-truck} 배송지 정하기 & 의사결정 트리 & 지도학습 \\
\hline
\end{tabular}
\end{center}
\end{frame}

% --- 1장: CNN ---
\begin{frame}{사과농장 1장: 수확 — 딥러닝 CNN}
\begin{block}{\emoji{deciduous-tree} 문제}
익은 사과와 덜 익은 사과를 구분해야 해요!
\end{block}
\vspace{0.3cm}
\begin{center}
\large CNN(합성곱 신경망) = 이미지의 특징을 자동으로 추출하여 분류
\end{center}
\vspace{0.3cm}
\begin{center}
\colorbox{lightblue}{\parbox{0.8\textwidth}{\centering
\emoji{light-bulb} CNN 시뮬레이터는 \textbf{다음 주(2주차 1차시)}에서 직접 체험합니다!\\
``사과 사진을 넣으면 AI가 어떻게 분류하는지 함께 확인해봐요''}}
\end{center}
\end{frame}

% --- 2장: KNN ---
\begin{frame}{사과농장 2장: 품종 분류 — KNN 시뮬레이터}
\begin{columns}
\begin{column}{0.48\textwidth}
\begin{center}
\small \textbf{학습 모드}\\[0.2cm]
\includegraphics[width=\textwidth]{images/fig13_KNN_학습모드.png}
\end{center}
\end{column}
\begin{column}{0.48\textwidth}
\begin{center}
\small \textbf{추론 모드}\\[0.2cm]
\includegraphics[width=\textwidth]{images/fig20_KNN_추론모드.png}
\end{center}
\end{column}
\end{columns}
\vspace{0.2cm}
\begin{center}
\small \emoji{light-bulb} ``가장 가까운 이웃이 많은 쪽으로 분류'' — 다수결 원리!
\end{center}
\end{frame}

% --- 4장: 선형 회귀 ---
\begin{frame}{사과농장 4장: 가격 책정 — 선형 회귀 시뮬레이터}
\begin{center}
\includegraphics[width=0.7\textwidth]{images/fig14_선형회귀_학습모드.png}
\end{center}
\vspace{0.1cm}
\begin{center}
\small \emoji{light-bulb} ``데이터 점들 사이의 추세선'' — 품질 점수 $\to$ 가격 예측!
\end{center}
\end{frame}

% --- 5장: 의사결정 트리 ---
\begin{frame}{사과농장 5장: 배송지 결정 — 의사결정 트리}
\begin{center}
\includegraphics[width=0.65\textwidth]{images/fig_decision_tree.png}
\end{center}
\end{frame}

\begin{frame}{의사결정 트리 시뮬레이터}
\begin{center}
\includegraphics[width=0.7\textwidth]{images/fig15_의사결정트리_학습분석모드.png}
\end{center}
\vspace{0.1cm}
\begin{center}
\small \emoji{light-bulb} 예/아니오 질문으로 분류 — 스무고개와 같은 원리!
\end{center}
\end{frame}

% --- 3장: K-Means ---
\begin{frame}{사과농장 3장: 등급 분류 — K-Means 시뮬레이터}
\begin{center}
\includegraphics[width=0.7\textwidth]{images/fig16_Kmeans_학습모드.png}
\end{center}
\vspace{0.1cm}
\begin{center}
\small \emoji{light-bulb} ``라벨 없이 AI가 스스로 그룹을 만든다'' — \textbf{비지도학습}!
\end{center}
\end{frame}

\begin{frame}{K-Means: 클러스터 품질은 어떻게 측정할까?}
\begin{center}
\includegraphics[width=0.7\textwidth]{images/fig21_Kmeans_Inertia.png}
\end{center}
\vspace{0.1cm}
\begin{center}
\small Inertia가 \textbf{낮을수록} 클러스터 품질이 좋음 (점들이 중심에 가까움)
\end{center}
\end{frame}

% --- 강화학습 시뮬레이터 ---
\begin{frame}{강화학습 시뮬레이터: Q-러닝}
\begin{center}
\includegraphics[width=0.65\textwidth]{images/fig19_Q러닝개념도.png}
\end{center}
\vspace{0.1cm}
\begin{center}
\small 상태(S) $\to$ 행동(A) $\to$ 보상(R) $\to$ 새 상태(S')
\end{center}
\end{frame}

\begin{frame}{스네이크 강화학습 시뮬레이터}
\begin{center}
\includegraphics[width=0.7\textwidth]{images/fig17_스네이크RL_전체화면.png}
\end{center}
\vspace{0.1cm}
\begin{center}
\small 무작위 이동 $\to$ 먹이 찾는 전략 습득! 시행착오로 학습하는 AI
\end{center}
\end{frame}

% --- 사과농장 에필로그 ---
\begin{frame}{사과농장 에필로그: 스마트 과수원 완성!}
\begin{center}
\begin{tabular}{|>{\centering\arraybackslash}p{0.8cm}|p{3.5cm}|>{\centering\arraybackslash}p{2.5cm}|>{\centering\arraybackslash}p{2cm}|p{3cm}|}
\hline
\rowcolor{tableheader}
\textbf{장} & \textbf{작업} & \textbf{ML 기술} & \textbf{학습 유형} & \textbf{핵심} \\
\hline
1 & \emoji{deciduous-tree} 수확 & CNN & 지도학습 & 이미지 분류 \\
\hline
2 & \emoji{magnifying-glass-tilted-left} 품종 & KNN & 지도학습 & 이웃 다수결 \\
\hline
3 & \emoji{package} 등급 & K-Means & 비지도학습 & 자동 그룹화 \\
\hline
4 & \emoji{money-bag} 가격 & 선형 회귀 & 지도학습 & 추세선 예측 \\
\hline
5 & \emoji{delivery-truck} 배송 & 의사결정 트리 & 지도학습 & 조건 분기 \\
\hline
\end{tabular}
\end{center}
\vspace{0.2cm}
\begin{center}
\small \emoji{star} 하나의 농장에서 \textbf{여러 ML 기술}이 함께 사용됩니다!
\end{center}
\end{frame}

%%%%%%%%%%%%%%%%%%%%%%%%%%%%%%%%%%%%%%%%%%%%%
\section{\emoji{orange-book} 전개 3: 머신러닝 학습 과정 4단계}
%%%%%%%%%%%%%%%%%%%%%%%%%%%%%%%%%%%%%%%%%%%%%

% --- 그림 10: 4단계 프로세스 ---
\begin{frame}{머신러닝 4단계 프로세스}
\begin{center}
\includegraphics[width=0.7\textwidth]{images/fig10_머신러닝프로세스.png}
\end{center}
\end{frame}

\begin{frame}{4단계와 엔트리 연결}
\begin{center}
\begin{tabular}{|>{\centering\arraybackslash}p{0.5cm}|>{\centering\arraybackslash}p{2.8cm}|p{4cm}|p{3.5cm}|}
\hline
\rowcolor{tableheader}
& \textbf{단계} & \textbf{설명} & \textbf{엔트리 연결} \\
\hline
1 & \textbf{데이터 수집} & 학습할 데이터 모으기 (양$\cdot$질$\cdot$다양성) & 웹캠 사진 촬영 \\
\hline
2 & \textbf{학습 (Training)} & 패턴 찾기 (에포크, 가중치) & ``모델 학습하기'' 버튼 \\
\hline
3 & \textbf{평가 (Testing)} & 새 데이터로 정확도 확인 & 정확도 \% 확인 \\
\hline
4 & \textbf{예측 (Prediction)} & 실전 적용 & 새 사진 분류 \\
\hline
\end{tabular}
\end{center}
\vspace{0.3cm}
\begin{center}
\colorbox{lightblue}{\parbox{0.8\textwidth}{\centering
\emoji{green-apple} \textbf{사과농장도 이 4단계}: 사과 사진 수집(1) $\to$ CNN 학습(2) $\to$ 정확도 확인(3) $\to$ 새 사과 판별(4)}}
\end{center}
\end{frame}

% --- 그림 11: 편향된 데이터 ---
\begin{frame}{좋은 데이터 = 좋은 AI}
\begin{center}
\includegraphics[width=0.6\textwidth]{images/fig11_편향된데이터vs좋은데이터.png}
\end{center}
\vspace{0.1cm}
\begin{center}
\small \emoji{warning} 빨간 사과만 학습하면, 초록 사과를 ``사과가 아니다''라고 판단!
\end{center}
\end{frame}

%%%%%%%%%%%%%%%%%%%%%%%%%%%%%%%%%%%%%%%%%%%%%
\section{\emoji{memo} 정리}
%%%%%%%%%%%%%%%%%%%%%%%%%%%%%%%%%%%%%%%%%%%%%

\begin{frame}{오늘 배운 핵심 개념 4가지}
\begin{block}{\emoji{pushpin} AI $\supset$ ML $\supset$ DL}
AI $>$ 기계학습 $>$ 딥러닝 — 포함 관계
\end{block}

\begin{block}{\emoji{pushpin} 기계학습 = 데이터에서 패턴을 학습}
전통적 프로그래밍: 사람이 규칙 작성 vs 기계학습: 컴퓨터가 패턴 발견
\end{block}

\begin{block}{\emoji{pushpin} 3가지 학습 유형}
지도(정답 있음) / 비지도(라벨 없음) / 강화(보상으로 학습)
\end{block}

\begin{block}{\emoji{pushpin} 4단계 학습 과정}
데이터 수집 $\to$ 학습 $\to$ 평가 $\to$ 예측
\end{block}
\end{frame}

\begin{frame}{엔트리 AI 모델과의 연결}
\begin{center}
\Large 3\textasciitilde 4주차 엔트리 이미지 분류 모델 = \textbf{지도학습}
\end{center}
\vspace{0.5cm}
\begin{center}
\begin{tabular}{rcl}
데이터 수집 & $\to$ & 웹캠으로 사진 촬영 \\
학습 & $\to$ & ``모델 학습하기'' 클릭 \\
평가 & $\to$ & 정확도 \% 확인 \\
예측 & $\to$ & 새 사진 분류 \\
\end{tabular}
\end{center}
\vspace{0.5cm}
\begin{center}
\small \emoji{light-bulb} 오늘 배운 원리가 엔트리 실습에서 \textbf{그대로} 적용됩니다!
\end{center}
\end{frame}

\begin{frame}{다음 차시 예고}
\begin{center}
\Large \textbf{2주차 1차시: 데이터와 딥러닝}
\end{center}
\vspace{0.5cm}
\begin{itemize}
    \item CNN 사과분류 시뮬레이터 \textbf{라이브 시연}
    \item 데이터가 AI 성능에 미치는 영향
    \item 신경망의 기본 구조
\end{itemize}
\vspace{0.5cm}
\begin{center}
\colorbox{lightblue}{\parbox{0.7\textwidth}{\centering
\emoji{memo} \textbf{다음 주 퀴즈}: 5지선다 10문항(0.5점) + 단답형 5문항(1점) = 10점}}
\end{center}
\end{frame}

\begin{frame}[standout]
\emoji{clapping-hands} 감사합니다!

\vspace{0.5cm}
질문이 있으신가요?
\end{frame}

\end{document}
